\documentclass[10pt]{article}
\usepackage[T1]{fontenc}
\usepackage[utf8]{inputenc}
\usepackage[french]{babel}
\usepackage{amsmath,amssymb,mathrsfs,mathpazo}
\usepackage{geometry}
\geometry{a4paper,margin=2cm}
\usepackage{graphicx,array,tabularx,tabularray,colortbl,multirow,multicol,booktabs,enumitem,pifont,ulem,fancybox,diagbox,cancel}
\usepackage{pstricks,pstricks-add,pst-plot,pst-tree,pst-node,pst-text,pst-func,pst-math,pst-bspline,pst-3dplot,pst-eucl,pst-solides3d}
\let\pscalculate\undefined
\usepackage{pst-bezier}
\usepackage{tikz,pgfplots,tkz-tab}
\usetikzlibrary{arrows,arrows.meta,patterns,calc,positioning,trees,decorations.markings,shapes.geometric,shapes.misc,shapes.arrows,intersections,decorations.pathmorphing,decorations.pathreplacing,backgrounds,shadows,babel}
\pgfplotsset{compat=1.18}
\usepgfplotslibrary{fillbetween}
\usepackage[np]{numprint}
\usepackage{listings,makecell,pas-tableur,twoopt}
\let\barre\undefined
\usepackage[pstricks]{tablvar}
\definecolor{darkgreen}{RGB}{0,100,0}
\providecommand{\red}{\color{red}}
\providecommand{\blue}{\color{blue}}
\providecommand{\green}{\color{green}}
\providecommand{\black}{\color{black}}
\providecommand{\white}{\color{white}}
\providecommand{\gray}{\color{gray}}
\providecommand{\cyan}{\color{cyan}}
\providecommand{\magenta}{\color{magenta}}
\providecommand{\yellow}{\color{yellow}}
\providecommand{\violet}{\color{violet}}
\providecommand{\orange}{\color{orange}}
\providecommand{\euro}{\texteuro}
\providecommand{\EUR}{\texteuro}
\providecommand{\ieme}{\textsuperscript{e}}
\providecommand{\degres}{\ensuremath{^\circ}}
\providecommand{\degree}{\ensuremath{^\circ}}
\providecommand{\vv}[1]{\overrightarrow{#1}}
\providecommand{\Coord}[1]{#1}
\providecommand{\decosix}{}
\providecommand{\starredbullet}{\ensuremath{\bullet}}
\pagestyle{empty}
\setlength{\parindent}{0pt}
\usepackage{adjustbox}\DeclareUnicodeCharacter{2236}{:}\DeclareUnicodeCharacter{2190}{\ensuremath{\leftarrow}}
\providecommand{\R}{\mathbb{R}}
\providecommand{\N}{\mathbb{N}}
\providecommand{\Z}{\mathbb{Z}}
\providecommand{\Q}{\mathbb{Q}}
\providecommand{\D}{\mathbb{D}}
\providecommand{\C}{\mathbb{C}}
\providecommand{\e}{{\mathrm{e}}}
\providecommand{\vect}[1]{\overrightarrow{#1}}
\providecommand{\vectt}[1]{\overrightarrow{\mathrm{#1}}}
\providecommand{\vvt}[1]{\overrightarrow{\mathrm{#1}}}
\providecommand{\np}[2][]{{#2~#1}}
\providecommand{\vv}[1]{{\overrightarrow{#1}}}
\renewcommand{\d}{{\mathrm{d}}}
\providecommand{\degres}{{^\circ}}
\providecommand{\degre}{{^\circ}}
\providecommand{\blue}{\color{blue}}
\providecommand{\red}{\color{red}}
\providecommand{\enskip}{\quad}
\providecommand{\llbracket}{\lbrack\!\lbrack}
\providecommand{\rrbracket}{\rbrack\!\rbrack}
\providecommand{\euro}{{\text{€}}}
\providecommand{\no}{{\text{n°}}}
\providecommand{\interval}[2]{{[#1;#2]}}
\providecommand{\intervFO}[2]{{[#1;#2[}}
\providecommand{\dotfill}{\ldots}

\providecommand{\ds}{\displaystyle}
\providecommand{\textding}[1]{#1}
\providecommand{\Oij}{\(\mathrm{O};\vec{\imath},\vec{\jmath}\)}
\providecommand{\Oijk}{\(\mathrm{O};\vec{\imath},\vec{\jmath},\vec{k}\)}
\providecommand{\Ouv}{\(\mathrm{O};\vec{u},\vec{v}\)}
\providecommand{\intervFF}[2]{\ensuremath{[#1;#2]}}
\renewcommand{\euro}{\ifmmode\text{\texteuro}\else\texteuro\fi}
\DeclareUnicodeCharacter{00B0}{\ifmmode^{\circ}\else\textdegree\fi}
\DeclareUnicodeCharacter{20AC}{\ifmmode\text{\texteuro}\else\texteuro\fi}
% Énoncés : Éducation nationale. Transcription et corrigés : APMEP, Denis Vergès et contributeurs. Découpage et classement : CoopMaths, Sébastien Lozano, Éric Elter et contributeurs. Adaptation de la présentation : Maths Explorer.
% CC BY-SA : https://coopmaths.fr/www/about/
% https://coopmaths.fr/alea/static/bac/2025/tex/bac_2025_06_sujet2_etrangers_4_cor.tex
\begin{document}
\noindent\textit{APMEP / CoopMaths. CC BY-SA. Présentation adaptée par Maths Explorer.}\par\medskip
\begin{enumerate}
\item 
	\begin{enumerate}
	\item $\vect{AB}\begin{pmatrix}-3\\1\\-1\end{pmatrix}$ et $\vect{AC}\begin{pmatrix}-1\\3\\-1\end{pmatrix}$. Ces deux vecteurs ne sont pas colinéaires car $\dfrac{-3}{-1}\neq \dfrac{1}{3}$. Les points $A$, $B$ et $C$ ne sont pas alignés.
	\item D'après la question précédente, les vecteurs $\vect{AB}$ et $\vect{AC}$ forment une base du plan $(ABC)$.
	Avec $\vect{n}\begin{pmatrix}-1\\1\\4\end{pmatrix}$, on calcule $\vect{n}\cdot\vect{AB}$ et $\vect{n}\cdot\vect{AC}$ :
	
	$\vect{n}.\vect{AB}=3+1-4=0$ et $\vect{n}.\vect{AC}=1+3-4=0$. $\vect{n}$ est donc orthogonal à deux vecteurs non colinéaires du plan $(ABC)$ donc il est normal à ce plan $(ABC)$.
	\item $\vect{n}$ étant normal au plan $(ABC)$, une équation cartésienne de ce plan est de la forme : $-x+y+4z+d=0$ avec $d \in \R$.

$A(1~;~0~;~3) \in (ABC)$ donc $-1+0+12+d=0 \iff d=-11$.

Donc $(ABC)$ : $-x+y+4x-11=0$.
	\end{enumerate}
\end{enumerate}

On considère le plan $\mathcal{P}$ d'équation cartésienne $3x - 3y + 2z - 9 = 0$ et le plan $\mathcal{P}'$ d'équation cartésienne $x - y- z + 2 = 0$.

\begin{enumerate}[resume]
	\item 
	\begin{enumerate}
		\item Le vecteur $\vect{n_1}\begin{pmatrix}3\\-3\\2\end{pmatrix}$ est normal au plan $\mathcal{P}$, le vecteur $\vect{n_2}\begin{pmatrix}1\\-1\\-1\end{pmatrix}$ est normal au plan $\mathcal{P}'$ . Les vecteurs $\vect{n_1}$ et $\vect{n_2}$ ne sont pas colinéaires car $\dfrac{3}{1}\neq \dfrac{2}{-1}$. Donc les plans $\mathcal{P}$ et $\mathcal{P}'$ ne sont ni parallèles ni confondus. Ils sont donc sécants. Leur intersection est donc une droite $(d)$.
		\item $\vect{n_1}\cdot \vect{n_2}=3+3-2=4\neq 0$ donc les plans ne sont pas perpendiculaires.
\end{enumerate}
	\item Soit une droite dirigée par le vecteur $\vect{u}\begin{pmatrix}1\\1\\0\end{pmatrix}$. Vérifions si $\vect{u}$ est normal aux vecteurs normaux des plans $\mathcal{P}$ et $\mathcal{P}'$.
	
	$\vect{n_1}.\vect{u}=3-3+0=0$ et $\vect{n_2}.\vect{u}=1-1+0=0$ donc cette droite est parallèle à la fois au plan $\mathcal{P}$ et au plan $\mathcal{P}'$. Donc il existe une droite $(d_1)$ appartenant à $\mathcal{P}$ et une droite $(d_1)$ appartenant à $\mathcal{P}'$ qui sont parallèles de vecteur directeur $\vect{n}$.
	Citons le théorème du toit : s'il existe deux droites parallèles $(d_1)$ et $(d_2)$, contenues respectivement dans deux plans $\mathcal{P}_1$ et $\mathcal{P}_2$, alors la droite, intersection des plans $\mathcal{P}_1$ et $\mathcal{P}_2$ est parallèle à $(d_1)$ et $(d_2)$.
	
Donc $(d)$ est parallèle à cette droite. Un vecteur directeur de $(d)$ est donc $\vect{n}$.
	\item On remplace les coordonnées de $M$ dans les deux équations de plans :
\begin{itemize}
\item $6-3+6-9=12-12=0$ donc $M\in \mathcal{P}$ ;
\item $2-1-3+2=4-4=0$ donc $M\in \mathcal{P}'$.
\end{itemize}

La droite $(d)$ passe donc par $M$ et a pour vecteur directeur $(d)~:~\vect{u}$ donc 

$\begin{cases} x=2+t\\y=1+t\\z=3\end{cases},~t\in \R$

	\item On cherche l'intersection de $(ABC)$ et de $(d)$. On résout le système : 

$\begin{cases} x=2+t\\y=1+t\\z=3\\-x+y+4z-11=0\end{cases}$. S'il admet une infinité de solutions alors $(d)\subset (ABC)$.

	On remplace $x$, $y$ et $z$ dans la dernière équation : 

	$-x+y+4z-11=0 \Rightarrow -(2+t)+1+t+4\times 3-11=0 \iff 0=0$. La dernière équation est vérifiée quelque soit la valeur de $t$. Donc ce système admet une infinité de solutions. $(d)\subset (ABC)$

	Donc la droite $(d)$ appartient aux trois plans $(ABC)$, $\mathcal{P}$ et $\mathcal{P}'$.
	Reprenons les trois vecteurs normaux des trois plans : $\vect{n}\begin{pmatrix}-1\\1\\4\end{pmatrix}$, $\vect{n_1}\begin{pmatrix}3\\-3\\2\end{pmatrix}$ et $\vect{n_2}\begin{pmatrix}1\\-1\\-1\end{pmatrix}$.
	
	$\vect{u}$ et $\vect{n_1}$ ne sont pas colinéaires, de même pour $\vect{u}$ et $\vect{n_2}$. Le plan $(ABC)$ n'est donc  confondu ni avec $\mathcal{P}$  ni avec $\mathcal{P}'$. Ces trois plans sont simplement sécants, d'intersections $(d)$.
\end{enumerate}

\end{document}
